\documentclass[12pt,a4paper]{article}
\usepackage{amsmath,amsfonts,amssymb}
\usepackage{graphicx}
\usepackage{hyperref}
\usepackage{enumitem}
\usepackage{geometry}
\geometry{margin=2.5cm}
\emergencystretch=2em

\begin{document}

\begin{center}
    {\Large \textbf{Master Thesis Project in Quantum Information Science and Technology}}\\[0.3cm]
    {\large \textbf{Quantum Machine Learning and Quantum Optimization for Finance: From Quantum Kernels and Neural Networks to Learned Initialization of the Quantum Approximate Optimization Algorithm}}\\[0.2cm]
    \vspace{0.3cm}
    {Duration: One Year (60 ECTS)}\\[0.2cm]
    {Student: \textit{NN}. Supervisor(s): \textit{Morten Hjorth-Jensen}, Department of Physics, University of Oslo}
\end{center}

\vspace{1cm}

\section*{1. Introduction}

Quantum machine learning (QML) combines parameterized quantum circuits
with classical learning algorithms~[1,2].  A circuit whose gates
depend on the data and on adjustable angles is run on a quantum
computer, a classical computer processes the measured expectation
values and updates the angles, and the two are trained together.  The
same hybrid structure is found in quantum kernel methods and quantum
neural networks~[3] and in the variational algorithms for
optimization, of which the quantum approximate optimization algorithm
(QAOA)~[4,5] is the most studied.

Finance is a natural test ground for these methods~[6], because it
offers both types of problems they are meant for.  The first is
learning from data: forecasting, classification and anomaly detection
for time series of asset prices, interest rates, volatilities and risk
factors.  The second is discrete optimization: the selection of a
portfolio with a fixed number of assets is a quadratic problem in
binary variables, of the same form as the graph problems on which QAOA
is normally tested.  For both types the open questions are what it
costs to train the quantum model and whether it does better than a
well-tuned classical method that is given the same data and the same
computing budget.

For QAOA the cost of training is the main obstacle.  Finding the
optimal angles of a variational circuit is NP-hard in general~[7],
gradients can vanish exponentially with the number of qubits~[8], and
every evaluation of the objective function requires many repetitions
of the circuit.  Two strategies reduce this cost.  A \emph{warm
start}~[9,10] replaces the uniform superposition by an initial state
that is biased toward a candidate solution from a classical algorithm.
\emph{Parameter initialization} supplies good angles before the
optimization starts, either by transfer from similar problem
instances~[11,12] or from a neural network trained on many
instances~[13].  The two strategies are normally treated separately.
A warm start, however, changes both the initial state and the mixer
Hamiltonian of the algorithm, and therefore the objective function in
which the angles are optimized.  A recent study~[14] proposes to learn
the two together: a graph neural network (GNN) with two output heads
predicts the warm-start state and the QAOA angles of a graph, and is
trained end to end through the expectation value of the quantum
circuit.  For small MaxCut instances the learned initializations gave
better cuts than a standard QAOA optimization and than a GNN that
solves MaxCut directly.  The joint model had only a slight advantage
over a network that predicts the angles alone, and the learned warm
starts stayed close to the uniform superposition.  Why this happens,
and whether it changes for larger graphs, deeper circuits, better
graph encoders and problems from finance, is not known.

European quantum computing infrastructure is at the same time becoming
available to academic projects.  The EuroHPC quantum computer VLQ of
the LUMI-Q consortium, a superconducting device with a star-shaped
qubit topology hosted by IT4Innovations in Ostrava, was inaugurated in
2025, and Norway takes part in the consortium through Sigma2.
Together with IBM Quantum and other cloud services, this makes it
realistic to run the hybrid algorithms of the thesis on quantum
hardware and not only on simulators.

This thesis project investigates classical and quantum methods for
machine learning and optimization problems that are relevant for
finance.  It has a common first part, in which the classical baselines
and a QAOA simulator are built, and three possible paths:
\begin{itemize}
  \item[\textbf{A.}] supervised learning for financial time series
  with quantum kernels and quantum neural networks;
  \item[\textbf{B.}] QAOA for MaxCut and for portfolio selection,
  including warm starts;
  \item[\textbf{C.}] neural networks that learn warm starts and QAOA
  angles jointly.
\end{itemize}
Selected algorithms are run on quantum hardware.  A thesis is expected
to carry out the common part and one or two of the paths in depth.
Paths~B and~C belong together, with B providing the tools and the
reference results for C.  The choice is made together with the
supervisors after the first months.

\section*{2. Problems, Methods and Scientific Motivation}

\textbf{Supervised learning with quantum circuits.}  A data point
$\boldsymbol x\in\mathbb R^d$ is encoded in the state
$|\phi(\boldsymbol x)\rangle=U(\boldsymbol x)|0\rangle^{\otimes n}$ of
$n$ qubits.  In a quantum support vector machine (QSVM) the quantum
computer estimates the kernel
\[
  k(\boldsymbol x,\boldsymbol x')
  =\bigl|\langle\phi(\boldsymbol x)|\phi(\boldsymbol x')\rangle\bigr|^2 ,
\]
which then enters a classical support vector machine~[3].  A quantum
neural network (QNN) adds a trainable circuit $W(\boldsymbol\theta)$
and uses an expectation value as model output,
\[
  f_{\boldsymbol\theta}(\boldsymbol x)
  =\langle\phi(\boldsymbol x)|\,W^\dagger(\boldsymbol\theta)\,O\,
  W(\boldsymbol\theta)\,|\phi(\boldsymbol x)\rangle ,
\]
with gradients obtained from the parameter-shift rule~[15].  Two
results set the frame for a comparison with classical models.  For
expressive feature maps the kernel values concentrate exponentially
around a constant as the number of qubits grows, and the number of
measurements needed to resolve them becomes prohibitive~[16].  For
classical data a quantum kernel can help only if its geometry differs
from that of the classical kernels in a way that matches the
labels~[17].

\textbf{The quantum approximate optimization algorithm.}  The
reference problem is weighted MaxCut on a graph $G=(V,E)$ with $n$
nodes,
\[
  \max_{\boldsymbol x\in\{0,1\}^n} C(\boldsymbol x),\qquad
  C(\boldsymbol x)=\sum_{(i,j)\in E}w_{ij}\,(x_i+x_j-2x_ix_j) ,
\]
which with $x_i=(1-Z_i)/2$ becomes the cost Hamiltonian
$H_C=\tfrac12\sum_{(i,j)\in E}w_{ij}(1-Z_iZ_j)$.  QAOA of depth $p$
prepares the state
\[
  |\boldsymbol\gamma,\boldsymbol\beta\rangle
  =\prod_{k=1}^{p}e^{-i\beta_kH_M}\,e^{-i\gamma_kH_C}\,|\psi_0\rangle ,
  \qquad H_M=-\sum_{i}X_i,\qquad |\psi_0\rangle=|+\rangle^{\otimes n},
\]
and a classical optimizer maximizes
$F_p(\boldsymbol\gamma,\boldsymbol\beta)=\langle\boldsymbol\gamma,\boldsymbol\beta|H_C|\boldsymbol\gamma,\boldsymbol\beta\rangle$
over the $2p$ angles.  The quality of the result is measured by the
approximation ratio $r=F_p/C_{\max}$.  The circuit can be read as a
Trotterized adiabatic evolution that starts in the ground state of
$H_M$ and has variationally chosen step lengths.  At $p=1$ it
guarantees $r\ge 0.6924$ on 3-regular graphs~[4], below the $0.878$ of
the classical Goemans--Williamson algorithm~[18], so the depth, the
initial state and the angles all matter.

\textbf{Portfolio selection.}  Choosing $B$ out of $n$ assets with
expected returns $\boldsymbol\mu$, covariance matrix $\Sigma$ and risk
aversion $q$ leads to the mean--variance problem~[19]
\[
  \min_{\boldsymbol x\in\{0,1\}^n}\;
  q\,\boldsymbol x^{T}\Sigma\,\boldsymbol x-\boldsymbol\mu^{T}\boldsymbol x ,
  \qquad \sum_{i=1}^{n}x_i=B .
\]
In qubit form this is an Ising Hamiltonian on a complete graph, with
couplings given by $\Sigma$ and local fields given by
$\boldsymbol\mu$~[20].  The budget constraint is included either as a
penalty $\lambda(\sum_ix_i-B)^2$ or exactly, with a mixer
$\sum_{i<j}(X_iX_j+Y_iY_j)$ that conserves the number of selected
assets and a Dicke state as initial state~[21].  Low-depth QAOA for
this problem has been benchmarked in Ref.~[22].

\textbf{Warm starts.}  Let $x_i\in[\varepsilon,1-\varepsilon]$ be a
candidate solution, for example from a continuous relaxation of the
problem.  Warm-start QAOA~[9] uses
\[
  |\psi_0\rangle=\bigotimes_{i=1}^{n}R_Y(\theta_i)|0\rangle,\qquad
  \theta_i=2\arcsin\sqrt{x_i},\qquad
  H_M^{\rm ws}=-\sum_{i}\bigl(\sin\theta_i\,X_i+\cos\theta_i\,Z_i\bigr),
\]
so that $|\psi_0\rangle$ is again the ground state of the mixer.  For
$x_i=1/2$ the standard algorithm is recovered.  For $x_i\to0$ or $1$
the initial state becomes a bit string and the mixer commutes with
$H_C$, and the circuit then only multiplies the candidate by a phase;
this is the reason for the regularization $\varepsilon$.  QAOA that
starts from a single good bit string makes little progress also with
the standard mixer~[23].  Useful warm starts lie between these limits.

\textbf{Learned initialization.}  A GNN~[24] computes an embedding
$\boldsymbol h_i$ for every node by several rounds of message passing
on the problem graph.  In the cascading model of Ref.~[14] a first
head maps each embedding to a warm-start value, and a second head maps
the pooled embeddings together with these values to the angles,
\[
  x_i=\sigma\bigl(f_{\rm ws}(\boldsymbol h_i)\bigr),\qquad
  (\boldsymbol\gamma,\boldsymbol\beta)
  =f_{\rm par}\bigl(\mathrm{Pool}_i\,[\boldsymbol h_i\oplus x_i]\bigr).
\]
The network parameters $\Theta$ are trained on a distribution of
graphs with the loss
\[
  \mathcal L(\Theta)=-\,\mathbb E_{G}\Bigl[
  F_p\bigl(\boldsymbol\gamma_\Theta(G),\boldsymbol\beta_\Theta(G);
  \boldsymbol x_\Theta(G)\bigr)\Bigr],
\]
whose gradient is propagated through the simulated circuit by
automatic differentiation.  No optimal solutions are needed as labels.
After training, one forward pass gives an initialization for a new
graph, which is used as it is or refined with a few optimizer steps.
A GNN trained in the same unsupervised way on a relaxed cost function
is also a purely classical heuristic for these problems~[25], and it
is one of the baselines.

Three types of quantities are used to compare the methods.
\begin{itemize}
  \item \textbf{Solution quality.}  For the optimization problems the
  approximation ratio of the expectation value and of the best sampled
  bit string, and the probability of sampling the optimum.  For the
  learning problems accuracy, F1 score or mean squared error on test
  data that lie later in time than the training data, the calibration
  of the predictions, and Sharpe-like measures where a trading rule
  is defined.
  \item \textbf{Quantum optimization budget.}  The number of circuit
  evaluations, multiplied by the number of repetitions of each, that
  is needed to reach a given quality.  This is the quantity which warm
  starts and learned initializations are meant to reduce.  The cost of
  training the neural network is reported separately, since it is paid
  once for a whole class of instances.
  \item \textbf{Resources and noise.}  The number of qubits, the
  number of two-qubit gates and the depth after compilation to a given
  device, and the loss of quality under noise.
\end{itemize}

The project is organized around four questions:
\begin{enumerate}
  \item For financial time series, do quantum kernels and quantum
  neural networks improve on well-tuned classical models in accuracy,
  in the amount of training data needed or in robustness, and which
  properties of the feature map decide this?
  \item How does low-depth QAOA perform for MaxCut and for portfolio
  selection compared with classical algorithms, and how do warm starts
  and constraint-preserving mixers change the result?
  \item Can a neural network that predicts the warm start and the
  angles together reduce the quantum optimization budget and improve
  the solutions?  Does it transfer to larger graphs, other graph
  families, larger depth and portfolio problems, and what does the
  warm-start head learn?
  \item Which of the differences found in simulation survive noise and
  execution on quantum hardware?
\end{enumerate}

\section*{3. Objectives of the Thesis}

\subsection*{Objective 1: Foundations, baselines and reproduction (common part)}
\begin{itemize}
  \item Study variational quantum algorithms, quantum kernels, QAOA
  and its variants, graph neural networks, and the two financial
  problems (prediction for time series and mean--variance portfolio
  selection).
  \item Implement QAOA for MaxCut as a state-vector simulation, check
  it against PennyLane or Qiskit, and reproduce the $p=1$ result for
  3-regular graphs and the dependence of the approximation ratio on
  $p$ for small random graphs.
  \item Build the classical baselines: support vector machines with
  linear, polynomial and RBF kernels and neural networks (fully
  connected, optionally recurrent) for the time-series tasks; exact
  solutions, the Goemans--Williamson algorithm and simulated annealing
  for the optimization problems.
  \item Select and preprocess the financial data sets, and fix the
  metrics and the protocols for training, validation and testing, with
  splits ordered in time.
\end{itemize}

\subsection*{Objective 2 (Path A): Quantum kernels and quantum neural networks for financial time series}
\begin{itemize}
  \item Implement quantum feature maps and QSVM classifiers on
  simulators, with different data encodings (angle encoding, amplitude
  encoding, data re-uploading).
  \item Characterize the kernels: concentration of the kernel values
  with the number of qubits, alignment with the labels, geometric
  difference from the classical kernels~[16,17], and the number of
  measurements needed.
  \item Implement variational classifiers and regressors and study
  their trainability (variance of the gradients) as a function of the
  number of qubits and the circuit depth.
  \item Compare with the classical baselines as functions of the size
  of the training set, and test the robustness to noise in the data
  and to limited data.
\end{itemize}

\subsection*{Objective 3 (Path B): QAOA for MaxCut and portfolio selection}
\begin{itemize}
  \item Study the dependence on the depth $p$ and on the classical
  optimizer (COBYLA, SPSA, gradient-based optimization with the
  parameter-shift rule), with a consistent count of the quantum
  optimization budget.  Test the concentration and transfer of optimal
  angles between instances~[11] and the interpolation strategy of
  Ref.~[12].
  \item Implement warm starts from the Goemans--Williamson relaxation
  for MaxCut and from the continuous relaxation of the portfolio
  problem~[9,10], and study the dependence on the regularization
  $\varepsilon$.  Map the objective function over the angles of the
  last layer for cold and warm starts on the same instance, as in
  Ref.~[14].
  \item Build portfolio instances from historical returns for up to
  about twenty assets.  Compare the penalty formulation with the
  constraint-preserving mixer, and the expectation value with the
  conditional value at risk~[26] as objective.  Compare with exact
  solutions and classical heuristics as functions of the risk
  aversion, the budget and the correlation structure of the assets.
\end{itemize}

\subsection*{Objective 4 (Path C): Jointly learned warm starts and QAOA angles}
\begin{itemize}
  \item Reproduce the GNN prediction of QAOA angles~[13] and the
  cascading two-head model of Ref.~[14] for unweighted
  Erd\H{o}s--R\'enyi graphs with 6 to 12 nodes at $p=2$, together with
  the baselines used there: QAOA optimized with COBYLA, a GNN that
  solves MaxCut directly, and a GNN that predicts the angles only.
  \item Carry out the ablations that separate the two effects:
  independent against cascading heads, a warm-start head alone, and
  fixed warm starts from a classical relaxation combined with learned
  angles.  Report the quality without any optimizer step and the
  number of circuit evaluations needed to reach a fixed approximation
  ratio.
  \item Test the hypothesis of Ref.~[14] that the extraction of graph
  information limits the model, and that this is why the learned warm
  starts stay close to $|+\rangle^{\otimes n}$.  Compare graph
  encoders (convolutional, isomorphism and attention networks),
  node features, pre-training of the warm-start head on relaxed
  classical solutions, and regularization terms in the loss.
  \item Analyze what the warm-start head learns: whether it fixes the
  nodes whose assignment is easy and keeps a superposition on the
  difficult ones, and how its output correlates with the
  Goemans--Williamson solution and with local properties of the graph.
  \item Test the transfer to graphs larger than those in the training
  set (up to the 20 to 24 qubits that can be simulated), to regular
  and weighted graphs, and to larger depth, where the angles are
  predicted either layer by layer or as a smooth schedule.
  \item Extend the model to portfolio selection: weighted complete
  graphs defined by the covariance matrix, with the expected returns
  as node features, trained on synthetic and historical instances.
\end{itemize}

\subsection*{Objective 5: Noise and quantum hardware}
\begin{itemize}
  \item Repeat selected calculations with noise models (depolarizing
  noise and device noise models) and determine how noise changes the
  kernel values, the QAOA expectation values and the ranking of the
  methods.
  \item Adapt the circuits to the connectivity and the native gates of
  the devices (star topology of VLQ, heavy-hex lattice of IBM
  Quantum), and compare the two-qubit gate counts of sparse MaxCut
  graphs and dense portfolio graphs.
  \item Run small instances on VLQ and/or IBM Quantum, depending on
  the allocations obtained: a quantum kernel or classifier for Path~A,
  and QAOA at $p=1$ and $p=2$ with cold, warm and learned
  initializations for Paths~B and~C.
  \item Apply error mitigation (readout calibration, zero-noise
  extrapolation) and compare with the simulations.
\end{itemize}

\subsection*{Objective 6: Possible extensions}
\begin{itemize}
  \item Boltzmann machines and quantum Boltzmann machines~[27] as
  generative models for financial data.
  \item Reinforcement-learning agents that choose the QAOA angles or
  the steps of the optimizer~[28], in connection with the
  reinforcement-learning projects of the group.
  \item Variants of QAOA (recursive and multi-angle QAOA), and angle
  schedules derived from a Trotterized annealing ramp as a
  physics-motivated initialization.
  \item Neural networks and quantum circuits that respect the
  symmetries of the graph.
  \item Other optimization problems from finance in binary quadratic
  form, for example index tracking.
\end{itemize}

\section*{4. Methods and Tools}

\begin{itemize}
  \item \textbf{Classical machine learning and optimization:} Python
  with NumPy, SciPy, Scikit-Learn and PyTorch or JAX; a GNN library
  (PyTorch Geometric or Jraph); NetworkX, a semidefinite-programming
  solver for the Goemans--Williamson algorithm and an exact solver for
  reference optima.
  \item \textbf{Quantum algorithms:} PennyLane and Qiskit for
  circuits, kernels and gradients, with state-vector and noisy
  simulators and access to hardware backends.  The model of Ref.~[14]
  was trained with a PennyLane--Catalyst and JAX pipeline on a single
  GPU in a few hours, so that Path~C can be carried out with the GPU
  resources available to the group.
  \item \textbf{Hardware:} VLQ through EuroHPC or national
  allocations, and IBM Quantum through its cloud service.
  \item \textbf{Data:} publicly available financial time series
  (equities, indices, interest rates) and, subject to access
  constraints, data provided through NBIM or project partners;
  synthetic graph and portfolio instances.
  \item \textbf{Useful background:} quantum mechanics, machine
  learning (FYS-STK4155) and quantum computing and quantum machine
  learning (FYS5419), with the lecture notes and codes of Ref.~[29].
\end{itemize}
All codes and results will be documented in Jupyter notebooks so that
the thesis can be reproduced.

\section*{5. Work Plan and Milestones (One-Year Timeline)}

\subsection*{Semester 1 (Months 1--6)}

\textbf{Months 1--2: Theory and reproduction}
\begin{itemize}
    \item Literature study; QAOA simulator for MaxCut and reproduction
    of the reference results; choice and preprocessing of the data.
\end{itemize}
\textbf{Deliverable:} Short report with the reproduced QAOA results
and a description of the data sets.

\bigskip

\textbf{Months 3--4: Classical baselines and choice of path}
\begin{itemize}
    \item Classical machine-learning and optimization baselines;
    metrics and protocols; first calculations along the paths and
    choice of the main path.
\end{itemize}
\textbf{Deliverable:} Validated classical pipeline and a plan for the
chosen path or paths.

\bigskip

\textbf{Months 5--6: Main path, first part}
\begin{itemize}
    \item Path~A: quantum kernels and QSVM on simulators, kernel
    diagnostics.
    \item Path~B: depth and optimizer study; warm starts for MaxCut.
    \item Path~C: reproduction of the GNN angle prediction and of the
    cascading model.
\end{itemize}
\textbf{Deliverable:} Working quantum models on simulators, with a
first comparison with the classical baselines.

\subsection*{Semester 2 (Months 7--12)}

\textbf{Months 7--8: Main path, second part}
\begin{itemize}
    \item Path~A: quantum neural networks; systematic comparison with
    the classical models.
    \item Path~B: portfolio selection with penalty and
    constraint-preserving formulations.
    \item Path~C: ablations, graph encoders, transfer to larger
    graphs and to portfolio instances.
\end{itemize}
\textbf{Deliverable:} Complete simulation results for the chosen path.

\bigskip

\textbf{Months 9--10: Noise and hardware}
\begin{itemize}
    \item Noisy simulations; compilation to the devices; runs on VLQ
    and/or IBM Quantum; error mitigation.
\end{itemize}
\textbf{Deliverable:} At least one of the algorithms executed on
quantum hardware and compared with the simulations.

\bigskip

\textbf{Month 11: Extensions and consolidation}
\begin{itemize}
    \item One of the extensions of Objective~6, depending on progress;
    final production runs and figures.
\end{itemize}

\bigskip

\textbf{Month 12: Writing and final submission}
\begin{itemize}
    \item Completion of the thesis manuscript (methods and theory
    chapters are written continuously during the year) and
    preparation of the final presentation.
\end{itemize}
\textbf{Deliverable:} Completed thesis and defense presentation.

\section*{6. Expected Outcomes}

\begin{itemize}
    \item A documented classical pipeline for the selected financial
    data sets and optimization problems, which serves as the reference
    for all quantum results.
    \item Implementations of the quantum methods of the chosen path:
    quantum kernels and quantum neural networks, QAOA with warm starts
    and constraint-preserving mixers, or the neural network for joint
    prediction of warm starts and angles.
    \item A quantitative comparison of classical and quantum methods
    in solution quality, quantum optimization budget, resources and
    sensitivity to noise.
    \item For Paths~B and~C, an assessment of when warm starts and
    learned initializations reduce the cost of QAOA, and of what the
    learned warm starts represent.
    \item At least one algorithm executed on quantum hardware.
    \item A master thesis and, depending on the results, a
    contribution to a scientific article.
\end{itemize}

\section*{7. References}

\begin{enumerate}
    \item M.~Schuld and F.~Petruccione, \textit{Supervised Learning with Quantum Computers}, Springer (2018); C.~Conti, \textit{Quantum Machine Learning}, Springer (2024); M.~Zhao \textit{et al.}, \textit{A Tutorial on Quantum Machine Learning and Quantum Neural Networks}, arXiv:2504.16131 (2025).
    \item M.~Cerezo \textit{et al.}, \textit{Variational quantum algorithms}, Nat.\ Rev.\ Phys.\ \textbf{3}, 625 (2021).
    \item V.~Havl\'i\v{c}ek \textit{et al.}, \textit{Supervised learning with quantum-enhanced feature spaces}, Nature \textbf{567}, 209 (2019); M.~Schuld and N.~Killoran, \textit{Quantum Machine Learning in Feature Hilbert Spaces}, Phys.\ Rev.\ Lett.\ \textbf{122}, 040504 (2019).
    \item E.~Farhi, J.~Goldstone and S.~Gutmann, \textit{A Quantum Approximate Optimization Algorithm}, arXiv:1411.4028 (2014).
    \item K.~Blekos \textit{et al.}, \textit{A review on Quantum Approximate Optimization Algorithm and its variants}, Phys.\ Rep.\ \textbf{1068}, 1 (2024).
    \item D.~Herman \textit{et al.}, \textit{Quantum computing for finance}, Nat.\ Rev.\ Phys.\ \textbf{5}, 450 (2023); R.~Or\'us, S.~Mugel and E.~Lizaso, \textit{Quantum computing for finance: Overview and prospects}, Rev.\ Phys.\ \textbf{4}, 100028 (2019).
    \item L.~Bittel and M.~Kliesch, \textit{Training Variational Quantum Algorithms Is NP-Hard}, Phys.\ Rev.\ Lett.\ \textbf{127}, 120502 (2021).
    \item J.~R.~McClean, S.~Boixo, V.~N.~Smelyanskiy, R.~Babbush and H.~Neven, \textit{Barren plateaus in quantum neural network training landscapes}, Nat.\ Commun.\ \textbf{9}, 4812 (2018).
    \item D.~J.~Egger, J.~Mare\v{c}ek and S.~Woerner, \textit{Warm-starting quantum optimization}, Quantum \textbf{5}, 479 (2021).
    \item R.~Tate, M.~Farhadi, C.~Herold, G.~Mohler and S.~Gupta, \textit{Bridging Classical and Quantum with SDP initialized warm-starts for QAOA}, ACM Trans.\ Quantum Comput.\ \textbf{4}, 9 (2023).
    \item F.~G.~S.~L.~Brand\~ao, M.~Broughton, E.~Farhi, S.~Gutmann and H.~Neven, \textit{For Fixed Control Parameters the Quantum Approximate Optimization Algorithm's Objective Function Value Concentrates for Typical Instances}, arXiv:1812.04170 (2018); A.~Galda, X.~Liu, D.~Lykov, Y.~Alexeev and I.~Safro, \textit{Transferability of optimal QAOA parameters between random graphs}, Proc.\ IEEE Int.\ Conf.\ Quantum Computing and Engineering (QCE), 171 (2021).
    \item L.~Zhou, S.-T.~Wang, S.~Choi, H.~Pichler and M.~D.~Lukin, \textit{Quantum Approximate Optimization Algorithm: Performance, Mechanism, and Implementation on Near-Term Devices}, Phys.\ Rev.\ X \textbf{10}, 021067 (2020).
    \item N.~Jain, B.~Coyle, E.~Kashefi and N.~Kumar, \textit{Graph neural network initialisation of quantum approximate optimisation}, Quantum \textbf{6}, 861 (2022).
    \item T.~Smedshaug, M.~Barani and P.~Crespo del Granado, \textit{Jointly Learning Warm Starts and Parameters for QAOA: A differentiable graph neural network framework for coupled QAOA initialization in MaxCut}, poster, IEEE Quantum Week (2026).
    \item K.~Mitarai, M.~Negoro, M.~Kitagawa and K.~Fujii, \textit{Quantum circuit learning}, Phys.\ Rev.\ A \textbf{98}, 032309 (2018); M.~Schuld, V.~Bergholm, C.~Gogolin, J.~Izaac and N.~Killoran, \textit{Evaluating analytic gradients on quantum hardware}, Phys.\ Rev.\ A \textbf{99}, 032331 (2019).
    \item S.~Thanasilp, S.~Wang, M.~Cerezo and Z.~Holmes, \textit{Exponential concentration in quantum kernel methods}, Nat.\ Commun.\ \textbf{15}, 5200 (2024).
    \item H.-Y.~Huang \textit{et al.}, \textit{Power of data in quantum machine learning}, Nat.\ Commun.\ \textbf{12}, 2631 (2021).
    \item M.~X.~Goemans and D.~P.~Williamson, \textit{Improved approximation algorithms for maximum cut and satisfiability problems using semidefinite programming}, J.\ ACM \textbf{42}, 1115 (1995).
    \item H.~Markowitz, \textit{Portfolio Selection}, J.\ Finance \textbf{7}, 77 (1952).
    \item A.~Lucas, \textit{Ising formulations of many NP problems}, Front.\ Phys.\ \textbf{2}, 5 (2014).
    \item S.~Hadfield \textit{et al.}, \textit{From the Quantum Approximate Optimization Algorithm to a Quantum Alternating Operator Ansatz}, Algorithms \textbf{12}, 34 (2019); Z.~Wang, N.~C.~Rubin, J.~M.~Dominy and E.~G.~Rieffel, \textit{XY mixers: Analytical and numerical results for the quantum alternating operator ansatz}, Phys.\ Rev.\ A \textbf{101}, 012320 (2020).
    \item S.~Brandhofer \textit{et al.}, \textit{Benchmarking the performance of portfolio optimization with QAOA}, Quantum Inf.\ Process.\ \textbf{22}, 25 (2023).
    \item M.~Cain, E.~Farhi, S.~Gutmann, D.~Ranard and E.~Tang, \textit{The QAOA gets stuck starting from a good classical string}, arXiv:2207.05089 (2022).
    \item T.~N.~Kipf and M.~Welling, \textit{Semi-Supervised Classification with Graph Convolutional Networks}, Proc.\ ICLR (2017); K.~Xu, W.~Hu, J.~Leskovec and S.~Jegelka, \textit{How Powerful are Graph Neural Networks?}, Proc.\ ICLR (2019).
    \item M.~J.~A.~Schuetz, J.~K.~Brubaker and H.~G.~Katzgraber, \textit{Combinatorial optimization with physics-inspired graph neural networks}, Nat.\ Mach.\ Intell.\ \textbf{4}, 367 (2022).
    \item P.~Kl.~Barkoutsos, G.~Nannicini, A.~Robert, I.~Tavernelli and S.~Woerner, \textit{Improving Variational Quantum Optimization using CVaR}, Quantum \textbf{4}, 256 (2020).
    \item M.~H.~Amin, E.~Andriyash, J.~Rolfe, B.~Kulchytskyy and R.~Melko, \textit{Quantum Boltzmann Machine}, Phys.\ Rev.\ X \textbf{8}, 021050 (2018).
    \item S.~Khairy, R.~Shaydulin, L.~Cincio, Y.~Alexeev and P.~Balaprakash, \textit{Learning to Optimize Variational Quantum Circuits to Solve Combinatorial Problems}, Proc.\ AAAI Conf.\ Artif.\ Intell.\ \textbf{34}, 2367 (2020); M.~M.~Wauters, E.~Panizon, G.~B.~Mbeng and G.~E.~Santoro, \textit{Reinforcement-learning-assisted quantum optimization}, Phys.\ Rev.\ Research \textbf{2}, 033446 (2020).
    \item M.~Hjorth-Jensen, \textit{Quantum Computing and Quantum Machine Learning}, lecture notes and codes,\\ \url{https://github.com/CompPhysics/QuantumComputingMachineLearning}.
\end{enumerate}

\end{document}
